\section{Research Methodology}
\label{sec:methodology}

\subsection{Two-Study Cross-Cultural Computational Architecture}
To overcome the pervasive social desirability bias that invalidates self-reported surveys on relationship misconduct, this study implements a **Two-Study Cross-Cultural Computational Mixed-Methods Architecture**. As depicted in Figure \ref{fig:method_pipeline}, the empirical strategy integrates: (1) large-scale naturalistic web-scraping of public first-person relationship dispute disclosures, treating online adjudication repositories strictly as an unobtrusive observational window into authentic interpersonal dynamics, (2) a **Split-Source Validation Framework** to prevent single-source algorithmic bias, (3) a human-validated bilingual Large Language Model (LLM) annotation pipeline, and (4) moderated mediation path modeling (Hayes' PROCESS Model 7 / SEM) evaluated across Western and non-Western cultural cohorts.

\begin{figure*}[htbp]
\centering
\fbox{
\begin{minipage}{0.95\textwidth}
\small
\textbf{Study 1 (English, Individualistic Context):} Reddit $r/\text{AmItheAsshole}$ Archives ($N = 10,000$ first-person moral judgment seeking narratives with crowd verdicts)\\[1ex]
\textbf{Study 2 (Thai, Collectivist Context):} Pantip Relationship Dispute Corpus ($N = 4,250$ systematically filtered narratives with moral adjudication queries)\\[1ex]
\textbf{Split-Source Measurement Layer:}
\begin{itemize}[leftmargin=*,noitemsep]
    \item \textit{Source A (Narrator Discursive Constructs):} SBERT Embeddings + Fine-Tuned RoBERTa Classifiers $\rightarrow$ $WEP, LRA, SSA$
    \item \textit{Source B (Behavioral Outcome \& Social Adjudication):} Documented Actions + Crowd-Sourced Moral Verdicts (`YTA`/`NTA` / communal backlash) $\rightarrow$ $EIR$
\end{itemize}
\textbf{Psychometric Validation Layer:} Standardized Bilingual LLM Annotation Protocol (decoding Temperature = 0.10, triangulated with crowd-adjudicated moral verdicts)\\[1ex]
\textbf{Statistical Layer:} Moderated Mediation Path Modeling (PROCESS Model 7) + Cross-Cultural Multi-Group Invariance (MICOM)
\end{minipage}
}
\caption{Two-Study Cross-Cultural Computational Research Architecture and Split-Source Validation Framework.}
\label{fig:method_pipeline}
\end{figure*}

\subsection{Dual Corpus Construction: Unobtrusive Observational Repositories}

\subsubsection{Study 1 Corpus (English): Reddit $r/\text{AmItheAsshole}$ ($N = 10,000$)}
We extracted archival first-person conflict narratives posted between January 2022 and September 2026 from Reddit's \textit{r/AmItheAsshole}. This repository provides an optimal observational window because authors explicitly solicit third-party moral adjudication from an impartial audience. Crucially, authors provide comprehensive chronological accounts detailing their relationship ideals, situational justifications, and concrete behavioral interactions. Only completed submissions with official community verdicts (\texttt{YTA} [You're the Asshole], \texttt{NTA} [Not the Asshole], \texttt{ESH} [Everyone Sucks Here], \texttt{NAH} [No Assholes Here]) were retained.

\subsubsection{Study 2 Corpus (Thai): Pantip Relationship Disputes ($N = 4,250$)}
To assess the cross-cultural generalizability of the model within a non-Western, collectivist setting, we constructed a targeted Thai corpus from Pantip (the primary social adjudication forum in Thailand). Rather than utilizing ad-hoc manual curation, we implemented a **systematic, rule-based computational ingestion pipeline**:
\begin{enumerate}[leftmargin=*]
    \item \textbf{First-Person Moral Judgment Queries:} Filtering for submissions where the title or opening paragraph explicitly seeks communal moral adjudication using standard phrases: \textit{``Am I wrong for...''} (Phid-Mai), \textit{``Am I selfish if...''} (Hen-Kae-Tua-Mai), \textit{``Is it too much to...''} (Koen-Pai-Mai), or \textit{``Is it fair that...''} (Fair-Mai).
    \item \textbf{Egalitarian Parity Cues:} Filtering for co-occurring parity terminology: \textit{``50/50''}, \textit{``split evenly''} (Han-Khrueng), \textit{``equality''} (Thao-Thiam), \textit{``fair''}, or \textit{``share expenses''} (Share-Kha-Chai-Chai).
    \item \textbf{Domain Strata:} Posts were sampled across four domestic and relational tags: Romantic Problems (\textit{Panha Khwam Rak}), Married Life (\textit{Chiwit Khu}), Family Dynamics (\textit{Khrop Khrua}), and Salary Earners (\textit{Manut Ngoen Duean}).
\end{enumerate}

\subsubsection{Data Sanitization and Ethical Safeguards}
All data curation adhered strictly to the Association of Internet Researchers (AoIR) ethical standards and institutional IRB public exemption criteria (`irb/`). All usernames, partner names, and geographic locations were scrubbed upon ingestion; user accounts were transformed into salted SHA-256 cryptographic hashes, guaranteeing complete non-reidentifiability.

\subsection{Split-Source Construct Operationalization \& 1--7 Likert Rubrics}
To eliminate common algorithm and method bias, the measurement model enforces a split-source design:

\subsubsection{1. Weaponized Egalitarian Posturing (\textit{WEP} --- Predictor $X$)}
Captures the frequency, density, and offensive deployment of 50/50 equality and parity demands. Operationalized through SBERT cosine similarity against an egalitarian semantic centroid and fine-tuned RoBERTa classification on a 1.0 to 7.0 continuous scale:
\begin{itemize}[leftmargin=*]
    \item \textit{WEP1 (Lexical Parity Density):} Frequency of explicit 50/50 and fairness markers.
    \item \textit{WEP2 (Semantic Centroid Alignment):} Embedding proximity to progressive relational discourse.
    \item \textit{WEP3 (Instrumental Leverage Claims):} Deploying parity standards to demand partner concessions.
    \item \textit{WEP4 (Relational Policing Vocabulary):} Accusing the partner of unequal labor, weaponized incompetence, or privilege.
\end{itemize}

\subsubsection{2. Legitimized Rational Accounts (\textit{LRA} --- Moderator $W$)}
Captures the deployment of plausible situational excuses, workload stress, or mental health boundaries (\textit{``I Have My Reasons''} / \textit{``	extit{I Have My Reasons''}}):
\begin{itemize}[leftmargin=*]
    \item \textit{LRA1 (Plausible Situational Excuses):} Citing acute occupational fatigue, commuting stress, or demanding schedules.
    \item \textit{LRA2 (Well-Being / Mental Health Exemption):} Citing emotional boundaries (``protecting my peace,'' sensory overload) to excuse non-reciprocity.
    \item \textit{LRA3 (Intellectualized Reasonableness Claims):} Explicitly framing oneself as the objective, logical partner while dismissing partner objections as emotional.
    \item \textit{LRA4 (Constraint-Based Special Standing):} Asserting personal circumstances that render standard 50/50 division inapplicable to oneself.
\end{itemize}

\subsubsection{3. Self-Serving Attribution \& Licensing (\textit{SSA} --- Mediator $M$)}
Captures the cognitive distortion that balances past virtue against present non-performance:
\begin{itemize}[leftmargin=*]
    \item \textit{SSA1 (Moral Credit Compensation):} Invoking past emotional or progressive support to excuse selfish conduct.
    \item \textit{SSA2 (Condemning Partner Character):} Magnifying partner flaws or perceived privilege to justify unilateral withdrawal.
    \item \textit{SSA3 (Rationalized Duty Exemption):} Asserting that personal reasons fully absolve the actor from relationship duties.
\end{itemize}

\subsubsection{4. Exploitative Interpersonal Relations (\textit{EIR} --- Outcome $Y$)}
Operationalized through a **Split-Source Triangulation**: combining documented behavioral actions within the narrative with **Crowd-Sourced Moral Verdicts**:
\begin{itemize}[leftmargin=*]
    \item \textit{EIR1 (Domestic Labor Shifting):} Complete unilateral refusal of domestic upkeep.
    \item \textit{EIR2 (Financial Parasitism):} Shifting living costs onto the partner while insisting on personal autonomy.
    \item \textit{EIR3 (Interpersonal Sanctions):} Weaponizing progressive terms to gaslight the partner or deploy stonewalling.
    \item \textit{EIR4 (Crowd-Adjudicated Moral Condemnation):} Triangulated with community accountability verdicts (Reddit `YTA` proportion and Thai communal accountability ratios).
\end{itemize}

\subsection{Bilingual LLM Annotation Protocol and Controlled Decoding}
To ensure psychometric rigor and replicability across languages, we developed a standardized bilingual annotation pipeline using three leading frontier model families: (1) OpenAI GPT-4o (\texttt{gpt-4o-2024-08-06}), (2) Anthropic Claude 3.5 Sonnet (\texttt{claude-3-5-sonnet-20241022}), and (3) Google Gemini 1.5 Pro (\texttt{gemini-1.5-pro-002}), with strictly controlled decoding parameter \textbf{Temperature = 0.10} to minimize stochastic variance.

Construct operationalization followed the standardized 1--7 Likert codebook (Section \ref{subsec:construct_operationalization}). To ensure scoring stability, each conflict narrative was processed through prompt ensembles evaluating narrator discursive patterns ($WEP, LRA, SSA$), achieving high cross-model consensus ($\kappa \ge 0.82$). Furthermore, to anchor the computational extractions in real-world ground truth without single-source algorithmic bias, the behavioral outcome ($EIR$) was directly triangulated with authentic crowd-sourced moral verdicts from the respective community adjudication archives (Reddit \texttt{YTA}/\texttt{NTA} verdicts and Thai communal accountability ratios).

\subsection{Statistical Protocol and Measurement Invariance}
The structured dataset was analyzed using moderated mediation path modeling (Hayes' PROCESS Model 7 / SEM) via 5,000-resample non-parametric bootstrapping. To verify cross-cultural generalizability between Study 1 and Study 2, we conducted Measurement Invariance of Composite Models (MICOM) testing:
\begin{enumerate}[leftmargin=*]
    \item \textbf{Configural Invariance:} Established identical indicator assignments across both corpora.
    \item \textbf{Compositional Invariance:} Evaluated via permutation testing ($c \ge 0.95, p > 0.05$).
    \item \textbf{Scalar / Equality Invariance:} Assessed through multi-group analysis (PLS-MGA) to test whether structural path coefficients differ significantly between Western and Thai dyads.
\end{enumerate}
